% ============================================================
%  Documentação de Projeto — Template LaTeX  (iteração 5)
%  Compile com pdfLaTeX no Overleaf
% ============================================================
\documentclass[12pt, a4paper]{article}
\usepackage[utf8]{inputenc}
\usepackage[T1]{fontenc}

% Margens: 1 polegada = 2.54cm todos os lados
% header distance from top of page = 720 twips = 0.5in
\usepackage[
  a4paper,
  top=2.54cm, bottom=2.54cm,
  left=2.54cm, right=2.54cm,
  headheight=15pt, headsep=0.46cm
]{geometry}

\usepackage{mathptmx}   % Times New Roman
\usepackage{xcolor}

% Cabeçalho: Times bold itálico 10pt, direita, sem linha
\usepackage{fancyhdr}
\pagestyle{fancy}
\fancyhf{}
\fancyhead[R]{\fontsize{10}{12}\selectfont\bfseries\itshape
  Documentação de Projeto para o Sistema \textless{}nome do sistema\textgreater{}
  \enspace Página \thepage}
\renewcommand{\headrulewidth}{0pt}
\fancypagestyle{capa}{\fancyhf{}\renewcommand{\headrulewidth}{0pt}}

% Parágrafo: sem indentação, sem parskip automático
\setlength{\parindent}{0pt}
\setlength{\parskip}{0pt}

% Corpo de texto: ragged right (NÃO justificado — igual ao original)
\raggedright

% Headings: H1=18pt bold (before=24pt after=12pt), H2=14pt bold (14/14)
\usepackage{titlesec}
\titleformat{\section}{\fontsize{18}{22}\selectfont\bfseries}{\thesection.}{1em}{}
\titlespacing*{\section}{0pt}{24pt}{12pt}
\titleformat{\subsection}{\fontsize{14}{17}\selectfont\bfseries}{\thesubsection}{1em}{}
\titlespacing*{\subsection}{0pt}{14pt}{14pt}
\titleformat{\subsubsection}{\normalsize\bfseries}{\thesubsubsection}{1em}{}
\titlespacing*{\subsubsection}{0pt}{10pt}{6pt}
\setcounter{secnumdepth}{3}

% Sumário: "Tabela de Conteúdo" 18pt bold
% Seções 11pt bold, subseções 11pt normal
% Indentação subseções = 360 twips = 0.635cm
% Sem dots; número de seção com ponto
\usepackage{tocloft}
\renewcommand{\contentsname}{Tabela de Conteúdo}
\renewcommand{\cfttoctitlefont}{\fontsize{18}{22}\selectfont\bfseries}
\renewcommand{\cftaftertoctitle}{\par}
\setlength{\cftbeforetoctitleskip}{6pt}
\setlength{\cftaftertoctitleskip}{10pt}

\renewcommand{\cftsecfont}{\fontsize{11}{13}\selectfont\bfseries}
\renewcommand{\cftsecpagefont}{\fontsize{11}{13}\selectfont\bfseries}
\setlength{\cftsecindent}{0pt}
\setlength{\cftsecnumwidth}{1.5em}
\setlength{\cftbeforesecskip}{3pt}
\renewcommand{\cftsecaftersnum}{.}
\renewcommand{\cftsecleader}{\hfill}
\renewcommand{\cftsecdotsep}{\cftnodots}

\renewcommand{\cftsubsecfont}{\fontsize{11}{13}\selectfont\normalfont}
\renewcommand{\cftsubsecpagefont}{\fontsize{11}{13}\selectfont\normalfont}
\setlength{\cftsubsecindent}{0.635cm}
\setlength{\cftsubsecnumwidth}{2em}
\setlength{\cftbeforesubsecskip}{3pt}
\renewcommand{\cftsubsecleader}{\hfill}
\renewcommand{\cftsubsecdotsep}{\cftnodots}

\usepackage{tabularx}
\usepackage{array}
\newcolumntype{L}[1]{>{\raggedright\arraybackslash}p{#1}}
\usepackage[hidelinks]{hyperref}

% ============================================================
\begin{document}

% ══════════════════════════════════════════════════════════════
%  CAPA — vbox com posições absolutas pixel-perfeitas
%  Medidas do original (px→pt @ 120dpi): 1px = 0.6pt
% ══════════════════════════════════════════════════════════════
\thispagestyle{capa}
\setcounter{page}{1}

% vbox ocupa toda a altura do texto; \vskip posiciona cada elemento
% \hbox to \textwidth{\hfil CONTENT} = alinhamento à direita
\noindent\vbox to \textheight{%
  % linha horizontal: começa 22.8pt abaixo da margem
  \vskip 22.8pt
  \hbox{\rule{\textwidth}{2.5pt}}%
  % "Documentação de Projeto" 32pt: gap=42.6pt depois da linha
  \vskip 42.6pt
  \hbox to \textwidth{\hfil
    {\fontsize{32}{38}\selectfont\bfseries Documentação de Projeto}}%
  % "para o sistema" 20pt: gap=40.8pt (corrigido pela profundidade da fonte 32pt)
  \vskip 40.8pt
  \hbox to \textwidth{\hfil
    {\fontsize{20}{24}\selectfont\bfseries para o sistema}}%
  % "<Nome do Sistema>" 32pt: gap=27.0pt
  \vskip 27.0pt
  \hbox to \textwidth{\hfil
    {\fontsize{32}{38}\selectfont\bfseries
      \textless\textcolor{red}{Nome do Sistema}\textgreater}}%
  % "Versão 1.0" 14pt: gap=43.2pt
  \vskip 43.2pt
  \hbox to \textwidth{\hfil
    {\fontsize{14}{17}\selectfont\bfseries Versão 1.0}}%
  % autores linha 1, 12pt: gap=34.8pt
  \vskip 34.8pt
  \hbox to \textwidth{\hfil
    {\fontsize{12}{14}\selectfont
      Projeto de sistema elaborado pelo(s) aluno(s) AAAA}}%
  % autores linha 2: gap=0.0pt
  \vskip 0.0pt
  \hbox to \textwidth{\hfil
    {\fontsize{12}{14}\selectfont
      como parte da disciplina \textbf{Projeto de Software}.}}%
  % "<data de criação>" 14pt: gap=44.4pt
  \vskip 44.4pt
  \hbox to \textwidth{\hfil
    {\fontsize{14}{17}\selectfont\bfseries
      \textless\textcolor{red}{data de criação}\textgreater}}%
  \vfil
}

\newpage

% ══════════════════════════════════════════════════════════════
%  SUMÁRIO + HISTÓRICO (página 2)
% ══════════════════════════════════════════════════════════════
\setcounter{page}{2}

\tableofcontents

% Espaço antes do Histórico (~2 blank lines of 12pt = 24pt + before H1 = 24pt)
\vspace{18pt}

% "Histórico de Revisões" estilo H1 manual
\vspace{24pt}
{\fontsize{18}{22}\selectfont\bfseries Histórico de Revisões}\par
\vspace{12pt}

\begin{tabularx}{\textwidth}{|L{3.2cm}|L{2.2cm}|X|L{2cm}|}
  \hline
  \textbf{Nome} & \textbf{Data} & \textbf{Razões para Mudança} & \textbf{Versão} \\
  \hline
  & & & \\[14pt]
  \hline
  & & & \\[14pt]
  \hline
\end{tabularx}

\newpage
\setcounter{page}{1}

% ══════════════════════════════════════════════════════════════
%  SEÇÕES
% ══════════════════════════════════════════════════════════════
\section{Introdução}

Este documento agrega: 1) a elaboração e revisão de modelos de domínio e
2) modelos de projeto para o sistema \textless\textcolor{red}{nome do sistema}\textgreater.
A referência principal para a descrição geral do problema, domínio e requisitos
do sistema é o documento de especificação que descreve a visão de domínio do sistema.

\section{Modelos de Usuário e Requisitos}

\subsection{Descrição de Atores}

Nesta subseção é apresentado descrição de cada um dos atores que interagem com o sistema.

\subsection{Modelo de Casos de Uso e Histórias de Usuários}

Nesta subseção é apresentado o diagrama de casos de uso do sistema. Para cada um deles,
utilize um ID que possa servir de referência no restante do documento.
Por exemplo UC-01 para o Caso de Uso 01.

\subsection{Diagrama de Sequência do Sistema e Contrato de Operações}

Nesta subseção é apresentado o diagrama de sequência do sistema de pelo menos, 3 Casos
de Uso ou Histórias de Usuário descritos na Seção 2.3.

\vspace{6pt}
Formato para cada contrato de operação

\vspace{4pt}
\begin{tabularx}{\textwidth}{|L{4cm}|X|}
  \hline
  \textbf{Contrato}             & \\[6pt] \hline
  \textbf{Operação}             & \\[6pt] \hline
  \textbf{Referências cruzadas} & \\[6pt] \hline
  \textbf{Pré-condições}        & \\[6pt] \hline
  \textbf{Pós-condições}        & \\[6pt] \hline
\end{tabularx}

\section{Modelos de Projeto}

\subsection{Arquitetura}

Pode ser descrita com um diagrama apropriado da UML ou \quad C4 Model

\subsection{Diagrama de Componentes e Implantação.}

Diagramas de componentes do sistema. Diagrama de implantação mostrando onde os
componentes estarão alocados para a execução.

\subsection{Diagrama de Classes}

Diagrama de classes do sistema.

\subsection{Diagramas de Sequência}

Diagramas de sequência para realização de casos de uso.

\subsection{Diagramas de Comunicação}

Diagramas de comunicação para realização de casos de uso.

\subsection{Diagramas de Estados}

Diagramas de estados do sistema.

\section{Modelos de Dados}

Deve-se apresentar os esquemas de banco de dados e as estratégias de mapeamento
entre as representações de objetos e não-objetos.

\end{document}
